\documentclass[10pt,a4paper]{article}
\usepackage[margin=1.7cm]{geometry}
\usepackage{amsmath,graphicx,array,hyperref}
\setlength{\parindent}{0pt}
\setlength{\parskip}{5pt}
\newcommand{\field}[1]{\texttt{\detokenize{#1}}}
\begin{document}
\begin{center}\Large\bfseries jetbns: inputs for NPC transport\end{center}
\textbf{Transfer only \field{npc_monte_carlo_inputs.zip}.}

It contains this PDF and the data file \field{npc_prebreakout_models.h5}.
The repository includes optional validator \field{read_npc_inputs.py}. It needs
NumPy and h5py, but neither jetbns, LaTeX, nor the original hydrodynamic files.
\begin{verbatim}
python -m pip install numpy h5py
python read_npc_inputs.py
python read_npc_inputs.py --model ana_reference
\end{verbatim}
The first command to the optional reader validates the file and lists model names.
The second prints all input scalars at the selected pre-breakout snapshot as JSON.
Use \field{ana_reference} for a simple first integration check. Each scenario
is a candidate for a transport calculation; no NPC efficiency is inferred here.

\section*{Data contract}
The root schema is \field{jetbns.npc-model-export.v3}. Version 2 used conical
propagation without cocoon coupling; these new trajectories supersede it.
The group \field{/models/<name>/} contains 201 time samples of each quantity.
The last sample is at
\[
t_* = t_{\rm start}+0.99(t_{\rm bo}-t_{\rm start}).
\]
This is a fixed physical sampling prescription, independent of the ODE step.
It is close to breakout, but not the immediately preceding solver step.
The exact breakout time and the remaining time to breakout are group attributes.
The separate \field{/prebreakout_summary/} gives one endpoint per model.

Every numeric series has \field{unit} and \field{frame} attributes. Read actual
arrays for a time-dependent calculation, or use the final row for a locally
steady shock. Time and radius are laboratory coordinates; densities and the
assumed field are upstream rest-frame quantities. The jet head is used as an
approximate downstream fluid velocity. Its speed is distinct from the shock speed.
\begin{verbatim}
import h5py
with h5py.File("npc_prebreakout_models.h5") as f:
    g = f["models/ana_reference"]
    n_p = g["proton_number_density_cm3"][-1]
    n_n = g["free_neutron_number_density_cm3"][-1]
    B_u = g["upstream_magnetic_field_g"][-1]
    Gamma_rel = g["relative_lorentz_factor"][-1]
    beta_s_u = g["shock_beta_upstream_frame"][-1]
    beta_s_d = g["shock_beta_downstream_frame"][-1]
    length_u = g["effective_path_length_upstream_cm"][-1]
\end{verbatim}

\section*{Model endpoints}
\begin{center}\footnotesize
\begin{tabular}{lrrrrrr}
Model & $\Gamma_{\rm rel}$ & $Y_e$ & $n_p$ [cm$^{-3}$] & $n_n$ [cm$^{-3}$]
& $\tau_{n:p}$ & $\tau_{p:n}$\\\hline
% MODEL_TABLE
\end{tabular}\end{center}
The two density columns have different physical meanings, explained on the next
page. A zero neutron excess is allowed and is not replaced by a numerical floor.

\newpage
\section*{Target densities and collision rates}
Let $n_b=\rho_u/m_p$ use upstream proper mass density. The prescription is
\begin{align}
X_{n,0}&=\max(0,1-2Y_e)\frac{2}{\pi}\arctan(M_n/M_{>r}),&
X_{n,\rm free}&=X_{n,0}e^{-t/\tau_n},\\
n_p&=f_{p,\rm free}Y_en_b,& n_n&=X_{n,\rm free}n_b.
\end{align}
The defaults are $M_n=10^{-4}M_\odot$, $\tau_n=900$ s, and $f_{p,\rm free}=1$.
This neutron-skin formula is equation (7) of the arXiv version of Metzger et al.
\cite{metzger}. The exterior mass includes the Lorentz factor on a lab-time
slice; for numerical ejecta it sums the individual ballistic contributions.
Its isotropic-equivalent polar mass convention is an additional approximation
when applying a spherical skin prescription.

\textbf{Composition limit.} $Y_e\rho/m_p$ counts proton content, including bound
protons. Treating it as a free-proton density is an explicit upper-proxy assumption,
controlled by \field{proton_free_fraction}. The neutron estimate is phenomenological,
not a network abundance. Neutron-decay daughters are not added to the proton proxy,
and $Y_e$ is held at its reconstructed value. These approximations need review
for long times, nuclear targets, or large neutron loading.

For a fast projectile the exported target encounter rates and local depths are
\begin{align}
\nu_{n:p}&=n_p\sigma_{pn}c,&\nu_{p:n}&=n_n\sigma_{pn}c,\\
\ell_u&=\Delta r/\Gamma_a,&
\tau_{n:p}&=n_p\sigma_{pn}\ell_u,\quad \tau_{p:n}=n_n\sigma_{pn}\ell_u.
\end{align}
Here $\sigma_{pn}=3\times10^{-26}$ cm$^2$ and $\Delta r=r$ by default.
The colon means \emph{projectile:target}; the long dataset names use
\field{neutron_to_proton_optical_depth} and
\field{proton_to_neutron_optical_depth}. These are encounter depths, not
charge-conversion probabilities. Transport must supply branching ratios and
any additional channels such as proton--proton interactions.

\textbf{Length convention.} $r/\Gamma_a$ is the adopted local dynamical length.
It is not an integrated exterior neutron/proton column or a general Lorentz
transformation of a stationary shell. If the Monte Carlo uses another geometry,
recompute its columns from the primitive densities and its own residence times.
At the endpoint, \field{snapshot_columns} also supplies the outward frozen-profile
estimates
\[
\tau^{\rm col}_{i:j}=\int_{r_*}^{r_{\rm out}}n_j(r,t_*)\sigma_{pn}
\Gamma_a(r,t_*)[1-\beta_a(r,t_*)],dr.
\]
These include the relative-motion factor for an outward ultrarelativistic
projectile. They can be much smaller than the local dynamical depths. Full
transport must still account for profile evolution along each particle path.

\section*{Shock, field, and approximate energy diagnostics}
\[
\Gamma_{\rm rel}=\Gamma_h\Gamma_a(1-\beta_h\beta_a),\quad
B_u=10^{13}\,\mathrm{G}\left(10^6\,\mathrm{cm}/r\right)^2,\quad
\omega_{p,1}=\frac{e B_u}{m_pc}.
\]
For the optional cold, unmagnetized $\hat\gamma=4/3$ jump closure,
\[
C=4\Gamma_{\rm rel}+3,\quad
u_{s,u}^2=\frac{(\Gamma_{\rm rel}-1)(4\Gamma_{\rm rel}/3+1)^2}
{2+8(\Gamma_{\rm rel}-1)/9},\quad u_{s,d}=u_{s,u}/C.
\]
Use $\beta_s=u_s/\sqrt{1+u_s^2}$. Advected downstream number densities may be
estimated as $Cn_u$ only if composition is unchanged across the shock. A transverse
frozen-in field would give $B_d=CB_u$; the magnetic geometry is a transport choice.
These closures do not resolve a radiation-mediated shock structure.

The retained diagnostic estimates follow Kashiyama et al. \cite{kashiyama}:
\[
k_BT_d=k_B(\rho_uc^2\Gamma_{\rm rel}^2/a)^{1/4},\qquad
\xi_b=\frac{eB_u}{\sigma_{pn}m_pc^2 n_b},\qquad
\gamma_{\max}\simeq\min\left(\frac{2m_ec^2}{6k_BT_d},\xi_b\right).
\]
\field{gyration_parameter}, \field{pn_optical_depth}, and the energy limits
retain the baryon-target reference. Use $\omega_{p,1}/\nu$ for a specified
target population. A formal energy limit below one signals failure of that
relativistic estimate; it is not a physical particle Lorentz factor. No energy
spectrum, injection normalization, pitch-angle distribution, conversion
probability, or acceleration efficiency is supplied by these tables.

\newpage
\section*{Numerical checks and model limits}
Each exported model passed a combined resolution comparison: time step
0.01 versus 0.005 s; breakout radial grid 192 versus 384; exterior-mass grid
256 versus 512; numerical extrapolation grid 256 versus 512; and native
outflow epochs versus native epochs plus midpoints. If needed, all grids are
refined further by factors of two; the final resolution is recorded per model.
The 48-point ellipsoidal cocoon-density average is tested independently against
192 points.
Endpoint relative changes must be below 5 percent for breakout and NPC fields,
and below 10 percent for the auxiliary cocoon energy, width, pressure and jet
area. Both tolerances are stored per model. An independent 1536-point opacity integral must satisfy
$|\tau\beta'_s-1|<0.05$. The measured errors are stored under
\field{convergence}; these check numerical stability, not the physical assumptions.

The initial numerical implementation resampled the short simulation history
onto a sparse grid extending to seconds. It could miss the mass ejection almost
entirely. Every recorded epoch is now retained. The previous numerical model
rankings and zero-neutron conclusions are superseded.

Analytical references use a mass-normalized homologous broken power law with
an exponential fast tail, $r=\beta ct$, and $\beta_{\rm limit}=0.9$. Numerical
profiles use recorded velocities and a generalized Gaussian with width
$\Delta\beta=0.05$ and exponent 2, preserving the smooth tail. The published
propagation study uses Bernoulli asymptotic velocities \cite{gutierrez}; this
export therefore is not a numerical reproduction of its figures. Measured
mass flux is extrapolated as $t^{-5/3}$ and launch velocity as $t^{-1/4}$ after
the file ends. Kernels are restricted to $0<\beta<1$ without renormalizing the
lost wings, retaining the existing model convention.

The jet and cocoon evolve together, including delayed cocoon pressure and
pressure-driven jet collimation. The subgroup \field{jet_cocoon} gives the
cocoon width, energy, pressure, lateral speed, and jet area/opening angle on
the same time grid. Breakout follows the legacy grey
snapshot criterion $\int\kappa\rho\,dr=1/\beta'_s$, $\kappa=0.16$ cm$^2$ g$^{-1}$.
This is not moving-medium frequency-dependent radiative transfer. Engine
duration, power, opening angle, sampling lag, and energy supplied by breakout
are recorded. Changing composition does not change this fixed-opacity
propagation solution. The models are controlled inputs for sensitivity studies;
an order-unity depth in one channel does not establish a successful NPC cycle.

\includegraphics[width=\linewidth]{npc_prebreakout_overview.png}

\newpage
\section*{Coupled propagation and example evolution}
Following equations (10)--(19) of Guti\'errez et al.\ \cite{gutierrez},
\[
\dot z_h=\beta_h c,\quad \dot r_c=\beta_{c,\perp}c,\quad
\dot E_c=\eta L_j(t_{\rm eng})(\beta_j-\beta_h),\quad
P_c=\frac{E_c(t-\Delta t_c)}{3V_c},
\]
\[
V_c=\frac{2\pi}{3}(z_h-z_0)r_c^2,\quad
\Delta t_c=\frac{\sqrt{3}(z_h-z_0)}{2c},\quad
\Sigma_j=\pi\tan^2\theta_0\min[z_h,(z_0+\hat z)/2]^2.
\]
The reconfinement height satisfies
$\hat z=z_0+(\beta_j-\beta_{\rm ej})\sqrt{L_j/(\pi\beta_jcP_c)}$.
The axial ellipsoidal density average follows the legacy approximation.
Lateral expansion uses relativistic speed addition. The finite injection
height is resolved numerically; the exported convergence checks include
cocoon energy, width, pressure and jet area. These are pre-breakout jet-head
inputs; the detached cocoon's subsequent shock and radiation are separate.

% EVOLUTION_PLOTS
\begin{thebibliography}{9}\small
\bibitem{kashiyama} K. Kashiyama, K. Murase, P. M\'esz\'aros (2013),
\href{https://arxiv.org/abs/1304.1945}{arXiv:1304.1945}.
\bibitem{metzger} B. Metzger et al. (2015),
\href{https://arxiv.org/abs/1409.0544}{arXiv:1409.0544}, Eq. (7).
\bibitem{gutierrez} E. Guti\'errez et al. (2025),
\href{https://arxiv.org/abs/2408.15973}{arXiv:2408.15973}.
\end{thebibliography}
\end{document}
